% Mathématiques 6e — cours V3
% Source LaTeX rendue côté Astro/KaTeX.

\section{Objectifs}

Dans ce chapitre, on apprend à :

\begin{itemize}
\item lire un horaire sur une horloge ou une frise ;
\item distinguer un \textbf{instant} d’une \textbf{durée} ;
\item convertir entre heures, minutes et secondes ;
\item calculer une durée écoulée entre deux horaires ;
\item déterminer un horaire de départ ou d’arrivée ;
\item additionner et soustraire des durées ;
\item passer d’une écriture heures-minutes à une écriture décimale d’heure ;
\item résoudre des problèmes de trajet, de planning et d’enchaînement d’activités.
\end{itemize}

\section{1. Instant et durée}

Un \textbf{instant} repère un moment précis.

Par exemple :

\[
14\ \text{h}\ 35\ \text{min}
\]

désigne un horaire.

Une \textbf{durée} mesure l’écart entre deux instants.

Par exemple :

\[
1\ \text{h}\ 20\ \text{min}
\]

peut représenter la durée d’un trajet.

Il ne faut pas confondre l’heure affichée par une horloge et le temps écoulé entre deux horaires.

\section{2. Lire une horloge}

\coursefigure{/images/mathematiques/6eme/temps-horloge.svg}{Horloge analogique avec une aiguille des heures et une aiguille des minutes.}{La petite aiguille indique l’heure et la grande aiguille permet de lire les minutes ; l’horaire étudié doit être interprété en tenant compte de la position des deux aiguilles.}

Sur une horloge analogique, un tour complet de la grande aiguille correspond à :

\[
60\ \text{min}.
\]

Un tour complet de la petite aiguille correspond à :

\[
12\ \text{h}.
\]

\section{3. Relations entre les unités}

Les unités de durée ne suivent pas uniquement un système décimal.

On utilise :

\[
1\ \text{h}=60\ \text{min},
\]

\[
1\ \text{min}=60\ \text{s}.
\]

Donc :

\[
1\ \text{h}=3\,600\ \text{s}.
\]

On utilise également :

\[
1\ \text{jour}=24\ \text{h}.
\]

Ces relations doivent être raisonnées et non traitées comme de simples déplacements de virgule.

\section{4. Convertir des heures en minutes}

Convertissons :

\[
2\ \text{h}\ 15\ \text{min}
\]

en minutes.

Les deux heures correspondent à :

\[
2\times60=120\ \text{min}.
\]

On ajoute :

\[
15\ \text{min}.
\]

Donc :

\[
\boxed{
2\ \text{h}\ 15\ \text{min}
=
135\ \text{min}
}.
\]

\section{5. Convertir des minutes en heures et minutes}

Convertissons :

\[
155\ \text{min}.
\]

On cherche combien d’heures complètes sont contenues dans cette durée.

Deux heures correspondent à :

\[
2\times60=120\ \text{min}.
\]

Il reste :

\[
155-120=35\ \text{min}.
\]

Donc :

\[
\boxed{
155\ \text{min}
=
2\ \text{h}\ 35\ \text{min}
}.
\]

\section{6. Convertir des minutes et secondes}

Pour :

\[
4\ \text{min}\ 20\ \text{s},
\]

les quatre minutes représentent :

\[
4\times60=240\ \text{s}.
\]

On ajoute :

\[
20\ \text{s}.
\]

Ainsi :

\[
\boxed{
4\ \text{min}\ 20\ \text{s}
=
260\ \text{s}
}.
\]

\section{7. Fraction d’heure}

Certaines durées usuelles correspondent à des fractions simples d’heure.

Une demi-heure vaut :

\[
\dfrac{1}{2}\ \text{h}=30\ \text{min}.
\]

Un quart d’heure vaut :

\[
\dfrac{1}{4}\ \text{h}=15\ \text{min}.
\]

Trois quarts d’heure valent :

\[
\dfrac{3}{4}\ \text{h}=45\ \text{min}.
\]

Ces repères sont utiles pour les conversions et le calcul mental.

\section{8. Écriture heures-minutes et écriture décimale}

L’écriture :

\[
1\ \text{h}\ 30\ \text{min}
\]

n’est pas égale à :

\[
1{,}30\ \text{h}.
\]

En effet :

\[
30\ \text{min}
=
\dfrac{30}{60}\ \text{h}
=
0{,}5\ \text{h}.
\]

Donc :

\[
\boxed{
1\ \text{h}\ 30\ \text{min}
=
1{,}5\ \text{h}
}.
\]

\section{9. Exemple développé 1 — conversion en heures décimales}

Convertissons :

\[
2\ \text{h}\ 24\ \text{min}
\]

en heures décimales.

Les minutes représentent :

\[
\dfrac{24}{60}\ \text{h}.
\]

On simplifie :

\[
\dfrac{24}{60}
=
\dfrac{2}{5}
=
0{,}4.
\]

Donc :

\[
\boxed{
2\ \text{h}\ 24\ \text{min}
=
2{,}4\ \text{h}
}.
\]

Le nombre après la virgule représente une \textbf{fraction d’heure}, pas un nombre de minutes.

\section{10. Calculer une durée entre deux horaires}

Pour calculer une durée, on peut avancer par étapes jusqu’à des horaires faciles.

\coursefigure{/images/mathematiques/6eme/temps-frise.svg}{Frise chronologique permettant de décomposer le temps entre un départ et une arrivée.}{Décomposer un trajet en plusieurs intervalles simples permet de calculer la durée totale sans traiter les horaires comme des nombres décimaux.}

Cette méthode est particulièrement efficace lorsque l’horaire traverse une heure entière.

\section{11. Exemple développé 2 — durée d’un trajet}

Un train part à :

\[
14\ \text{h}\ 38\ \text{min}
\]

et arrive à :

\[
16\ \text{h}\ 05\ \text{min}.
\]

De :

\[
14\ \text{h}\ 38\ \text{min}
\]

à :

\[
15\ \text{h},
\]

il s’écoule :

\[
22\ \text{min}.
\]

De :

\[
15\ \text{h}
\]

à :

\[
16\ \text{h},
\]

il s’écoule :

\[
1\ \text{h}.
\]

Enfin, de :

\[
16\ \text{h}
\]

à :

\[
16\ \text{h}\ 05\ \text{min},
\]

il s’écoule :

\[
5\ \text{min}.
\]

La durée totale vaut :

\[
22\ \text{min}
+
1\ \text{h}
+
5\ \text{min}
=
1\ \text{h}\ 27\ \text{min}.
\]

\section{12. Soustraire deux horaires}

On peut aussi effectuer une soustraction organisée.

Pour :

\[
16\ \text{h}\ 05\ \text{min}
-
14\ \text{h}\ 38\ \text{min},
\]

on ne peut pas directement faire :

\[
5-38.
\]

On transforme :

\[
16\ \text{h}\ 05\ \text{min}
=
15\ \text{h}\ 65\ \text{min}.
\]

Alors :

\[
65-38=27\ \text{min},
\]

et :

\[
15-14=1\ \text{h}.
\]

On retrouve :

\[
1\ \text{h}\ 27\ \text{min}.
\]

\section{13. Additionner des durées}

Considérons :

\[
2\ \text{h}\ 48\ \text{min}
+
1\ \text{h}\ 37\ \text{min}.
\]

On additionne séparément les heures et les minutes :

\[
2+1=3\ \text{h},
\]

\[
48+37=85\ \text{min}.
\]

Or :

\[
85\ \text{min}
=
1\ \text{h}\ 25\ \text{min}.
\]

Donc :

\[
3\ \text{h}
+
1\ \text{h}\ 25\ \text{min}
=
\boxed{4\ \text{h}\ 25\ \text{min}}.
\]

\section{14. Soustraire des durées}

Calculons :

\[
3\ \text{h}\ 10\ \text{min}
-
1\ \text{h}\ 45\ \text{min}.
\]

On transforme :

\[
3\ \text{h}\ 10\ \text{min}
=
2\ \text{h}\ 70\ \text{min}.
\]

Alors :

\[
70-45=25\ \text{min},
\]

et :

\[
2-1=1\ \text{h}.
\]

Donc :

\[
\boxed{
3\ \text{h}\ 10\ \text{min}
-
1\ \text{h}\ 45\ \text{min}
=
1\ \text{h}\ 25\ \text{min}
}.
\]

\section{15. Déterminer un horaire d’arrivée}

Une activité commence à :

\[
9\ \text{h}\ 25\ \text{min}
\]

et dure :

\[
1\ \text{h}\ 55\ \text{min}.
\]

On ajoute :

\[
9\ \text{h}\ 25\ \text{min}
+
1\ \text{h}\ 55\ \text{min}.
\]

Cela donne :

\[
10\ \text{h}\ 80\ \text{min}.
\]

Comme :

\[
80\ \text{min}
=
1\ \text{h}\ 20\ \text{min},
\]

l’horaire final est :

\[
\boxed{
11\ \text{h}\ 20\ \text{min}
}.
\]

\section{16. Déterminer un horaire de départ}

Un trajet se termine à :

\[
18\ \text{h}\ 10\ \text{min}
\]

et dure :

\[
2\ \text{h}\ 35\ \text{min}.
\]

On peut remonter :

\[
18\ \text{h}\ 10\ \text{min}
-
2\ \text{h}\ 35\ \text{min}.
\]

On transforme :

\[
18\ \text{h}\ 10\ \text{min}
=
17\ \text{h}\ 70\ \text{min}.
\]

Donc :

\[
70-35=35\ \text{min},
\]

et :

\[
17-2=15\ \text{h}.
\]

Le départ a lieu à :

\[
\boxed{
15\ \text{h}\ 35\ \text{min}
}.
\]

\section{17. Exemple développé 3 — planning à plusieurs étapes}

Une séance commence à :

\[
9\ \text{h}\ 25\ \text{min}.
\]

Une première activité dure :

\[
1\ \text{h}\ 55\ \text{min}.
\]

Elle se termine à :

\[
11\ \text{h}\ 20\ \text{min}.
\]

Une pause de :

\[
35\ \text{min}
\]

mène à :

\[
11\ \text{h}\ 55\ \text{min}.
\]

Une seconde activité dure :

\[
50\ \text{min}.
\]

L’horaire final est donc :

\[
\boxed{
12\ \text{h}\ 45\ \text{min}
}.
\]

Dans un problème complexe, une frise ou un tableau d’étapes aide à éviter les oublis.

\section{18. Durée traversant minuit}

Une activité commence à :

\[
23\ \text{h}\ 35\ \text{min}
\]

et se termine le lendemain à :

\[
0\ \text{h}\ 20\ \text{min}.
\]

De :

\[
23\ \text{h}\ 35\ \text{min}
\]

à minuit, il reste :

\[
25\ \text{min}.
\]

Puis :

\[
20\ \text{min}
\]

s’écoulent après minuit.

La durée totale est :

\[
25+20=45\ \text{min}.
\]

Le changement de jour doit donc être pris en compte.

\section{19. Contrôler une durée}

Un ordre de grandeur peut permettre de détecter une erreur.

Si un trajet commence vers :

\[
14\ \text{h}\ 30
\]

et se termine vers :

\[
16\ \text{h},
\]

la durée est de l’ordre de :

\[
1\ \text{h}\ 30\ \text{min}.
\]

Un résultat de :

\[
15\ \text{min}
\]

ou de :

\[
15\ \text{h}
\]

serait manifestement incohérent.

\section{20. Méthode — convertir une durée}

\subsection{Étape 1 — identifier l’unité de départ}

Heures, minutes ou secondes.

\subsection{Étape 2 — utiliser la relation appropriée}

Par exemple :

\[
1\ \text{h}=60\ \text{min}.
\]

\subsection{Étape 3 — convertir chaque partie}

Une durée mixte peut nécessiter plusieurs calculs.

\subsection{Étape 4 — contrôler}

Une petite unité donne généralement une valeur numérique plus grande.

\section{21. Méthode — calculer un temps écoulé}

\subsection{Étape 1 — placer départ et arrivée}

Une frise peut aider.

\subsection{Étape 2 — choisir une stratégie}

\begin{itemize}
\item avancer jusqu’à des horaires simples ;
\item effectuer une soustraction avec regroupement ;
\item convertir tout en minutes lorsque cela simplifie réellement.
\end{itemize}

\subsection{Étape 3 — normaliser la réponse}

Une écriture finale comme :

\[
1\ \text{h}\ 85\ \text{min}
\]

doit être transformée.

\subsection{Étape 4 — interpréter}

Préciser qu’il s’agit d’une durée.

\section{22. Erreurs fréquentes}

\subsection{Confondre base dix et base soixante}

Il est faux d’écrire :

\[
1\ \text{h}\ 45\ \text{min}
=
1{,}45\ \text{h}.
\]

En réalité :

\[
45\ \text{min}
=
\dfrac{45}{60}\ \text{h}
=
0{,}75\ \text{h}.
\]

Donc :

\[
1\ \text{h}\ 45\ \text{min}
=
1{,}75\ \text{h}.
\]

\subsection{Laisser plus de soixante minutes dans une écriture finale}

Une écriture :

\[
2\ \text{h}\ 75\ \text{min}
\]

se simplifie en :

\[
3\ \text{h}\ 15\ \text{min}.
\]

\subsection{Soustraire directement les minutes sans regroupement}

Lorsque le nombre de minutes de l’horaire final est plus petit que celui de l’horaire initial, il faut transformer une heure en :

\[
60\ \text{min}.
\]

\subsection{Oublier le changement de jour}

Un horaire après minuit peut appartenir au lendemain.

\subsection{Confondre instant et durée}

Un départ à :

\[
8\ \text{h}\ 30
\]

n’est pas une durée de :

\[
8\ \text{h}\ 30.
\]

\section{23. À retenir}

Les relations fondamentales sont :

\[
\boxed{
1\ \text{h}=60\ \text{min}
}
\]

et :

\[
\boxed{
1\ \text{min}=60\ \text{s}
}.
\]

Une durée est l’écart entre deux instants.

Les calculs de temps utilisent fréquemment des regroupements par :

\[
60.
\]

Une écriture heures-minutes n’est pas une écriture décimale.

Par exemple :

\[
\boxed{
1\ \text{h}\ 30\ \text{min}=1{,}5\ \text{h}
}.
\]

Pour résoudre un problème :

\begin{itemize}
\item repérer départ, arrivée et durées intermédiaires ;
\item utiliser une frise si nécessaire ;
\item convertir dans une unité adaptée ;
\item gérer correctement les regroupements ;
\item contrôler l’ordre de grandeur ;
\item rédiger l’horaire ou la durée avec son unité.
\end{itemize}
